% =====================================================================
%  Section 1: Introduction
% =====================================================================
\section{Introduction}
\label{sec:intro}

\paragraph{Data selection in instruction tuning.}
Instruction tuning (IT) is the foundational stage of LLM
post-training: contemporary open releases~\cite{dubey2024llama3,
lambert2024tulu3,qwen2025} begin with supervised fine-tuning on
instruction--response pairs. As IT corpora have grown from
thousands to hundreds of thousands of pairs, an empirical
near-consensus has emerged: at fixed compute, the quality and
composition of the IT subset dominate its raw
scale~\cite{zhou2023lima,zhang2025jair, qin2024datatsunami}. The
central lever of IT is therefore \emph{selection}---choosing which
subset of a large candidate pool to train on.

\paragraph{Selection methods.}
Selection criteria have evolved along three model-aware
axes---surface-level quality scoring, gradient influence on a
target validation set, and structural signals from hidden-state
geometry (Sec.~\ref{sec:related}). A more recent line is
\emph{dynamic data selection}: re-estimating per-candidate utility
during training as the model state changes~\cite{kung2023activeit,
sun2025dots}. This line shows that useful examples are not fixed
properties of the data alone; their value can depend on the current
model, training stage, and optimisation regime. In this paper we
focus on the supervised instruction-tuning setting, where
candidates are instruction--response pairs and no rollout group,
verifier reward, or target validation set is assumed.

\paragraph{Two gaps in dynamic selection.}
Despite this progress, two gaps remain. First, structural
hidden-state selectors and dynamic selectors have largely evolved
separately. The former extract capability directions from model
representations but usually fix them before training using seed or
target examples. The latter adapt during training but typically
optimise scalar utilities such as difficulty, uncertainty, loss
change, or reward proxies without directly using the model's
evolving representation geometry. Dynamic selection can track the
training state while still missing the structural direction along
which the model is being updated.

Second, the scalar utilities used in dynamic selection can be
unstable in multi-task instruction pools. In the composite-reward
selector we adopt as a base, difficulty and uncertainty are
combined through a variance-ratio weight. In multi-task pools this
weight can shift with task composition because it conflates
within-task and between-task variance; under a naive
per-micro-batch implementation it can even lose the difficulty
signal entirely when the micro-batch size is one. These
sensitivities motivate a dynamic selection criterion that is not
only adaptive but also representation-aware and stable.

\paragraph{Our approach: Trajectory anchoring.}
We propose \textsc{TADS} (Trajectory-Anchored Dynamic Selection)
whose central component is a \emph{trajectory anchor}
$v_l^{(t)}$: a per-layer capability direction extracted directly
from the model's own hidden-state trajectory at each refresh step.
Concretely, at the start of epoch~$t$ we form contextualisation
deltas
\[
\Delta h_l(x;\theta_{t-1})
=
h_l^{\mathrm{last}}(x;\theta_{t-1})
-
h_l^{\mathrm{first}}(x;\theta_{t-1})
\]
on a probe subset of the candidate pool, and take the top principal
direction of their covariance as the layer-$l$ trajectory anchor
$v_l^{(t)}$. The anchor is model-intrinsic: it requires no
in-domain seed examples, no validation set, no labels, and no
auxiliary scorer. It is extracted from the very model being
fine-tuned, using the candidate pool itself. This anchor is the
key contribution of the paper; the remaining components of the
selection rule are built around it.

To integrate the anchor with the existing composite-reward signal
$\mathcal{R}_i^{(t)}$ used by prior dynamic
selectors~\cite{kung2023activeit, sun2025dots}, \textsc{TADS} forms
the per-candidate score as the product of a bounded calibrated
utility (derived from $\mathcal{R}_i^{(t)}$) and a
trajectory-anchor factor:
\[
s_i^{(t)}
=
\underbrace{\tilde{\mathcal{R}}_i^{(t)}}_{\substack{\text{calibrated utility}\\[1pt]\text{(supporting)}}}
\;\cdot\;
\underbrace{\bigl(1+\lambda\,\widetilde{\mathrm{align}}_i^{(t)}\bigr)}_{\substack{\text{trajectory-anchor factor}\\[1pt]\text{(central novelty)}}},
\]
where $\widetilde{\mathrm{align}}_i^{(t)}\in[0,1]$ is the
projection of the candidate's sequence-mean hidden state onto
$v_l^{(t)}$, averaged across layers. The calibrated utility is a
deterministic pool-level transform of $\mathcal{R}_i^{(t)}$
($\tilde{\mathcal{R}}_i^{(t)} = \sigma((\mathcal{R}_i^{(t)} -
\bar{\mathcal{R}}^{(t)})/(\sigma_{\mathcal{R}}^{(t)} +
\varepsilon)) \in (0,1)$) that simply bounds and de-biases the
composite reward; it serves as a stable carrier for the anchor.
Setting $\lambda=0$ disables the anchor and recovers the
calibrated-utility base ranking, yielding a clean ablation that
isolates the anchor's contribution.

\paragraph{Stability of the trajectory anchor.}
A natural question is whether the anchor is a meaningful structural
object or merely probe noise. We answer this with a stability
guarantee. Under a positive eigenvalue gap and bounded covariance
change, the sign-calibrated trajectory anchor satisfies
\[
\left\|
v_l^{(t+1)}-v_l^{(t)}
\right\|
\le
\frac{2C_\Sigma}{\gamma_{\min}}\,\eta_t .
\]
Thus the anchor evolves smoothly at the scale of the learning-rate
schedule. The guarantee concerns the trajectory anchor itself and
therefore the structural component introduced by \textsc{TADS}; it
does not depend on the particular form of the utility transform.

\paragraph{Contributions.}
This paper makes four contributions, organised around the
trajectory anchor. \emph{First and centrally}, we introduce the
\emph{trajectory anchor}: a per-layer capability direction
extracted online from the model's own contextualisation deltas at
each refresh step, with no in-domain seeds, labels, validation
set, or auxiliary scorer. The anchor is the structural object
TADS adds to the dynamic-selection toolkit. \emph{Second}, we
prove the anchor is stable: under a positive eigenvalue gap, its
epoch-to-epoch drift is bounded by the learning rate
($\|v_l^{(t+1)}-v_l^{(t)}\| \le (2C_\Sigma/\gamma_{\min})\eta_t$),
yielding what is to our knowledge the first explicit stability
guarantee for a dynamic IT data-selection criterion; the
guarantee concerns the structural anchor itself, not any
particular utility transform layered on top. \emph{Third}, we
identify a practical instability in composite-reward selection
that the anchor compensates for: the variance-ratio weight can
drift with task mix and can collapse silently under naive
micro-batch statistics. This motivates a representation-level
signal that scalar composite rewards alone cannot express.
\emph{Fourth}, we empirically evaluate \textsc{TADS} across model
scales, instruction-tuning datasets, and benchmark families,
showing improvements over strong static, seed-based, and dynamic
baselines with less than $0.5\%$ wall-clock overhead, while
$\lambda$ ablations isolate the contribution of the trajectory
anchor.
